% Ad Atticum 15.14 (ad-atticum-15-14)
% Source: https://raw.githubusercontent.com/PerseusDL/canonical-latinLit/master/data/phi0474/phi057/phi0474.phi057.perseus-lat1.xml
% Fetched by scripts/fetch_latin.py

% Dateline (Perseus): Scr. in Tusculano v K. Quint. a. 710 (44).

\ciceroLetterOpener{CICERO ATTICO salutem}

\ciceroSection{1} vi Kalend. accepi a Dolabella litteras. quarum exemplum tibi misi. in quibus erat omnia se fecisse quae tu velles. statim ei rescripsi et multis verbis gratias egi. sed tamen ne miraretur cur idem iterum facerem, hoc causae sumpsi quod ex te ipso coram antea nihil potuissem cognoscere. sed quid multa? litteras hoc exemplo dedi:

\ciceroSection{2} CICERO DOLABELLAE COS. suo. antea cum litteris Attici nostri de tua summa liberalitate summoque erga se beneficio certior factus essem cumque tu ipse etiam ad me scripsisses te fecisse ea quae nos voluissemus, egi tibi gratias per litteras iis verbis ut intellegeres nihil te mihi gratius facere potuisse. postea vero quam ipse Atticus ad me venit in Tusculanum huius unius rei causa tibi ut apud me gratias ageret, cuius eximiam quandam et admirabilem in causa Buthrotia voluntatem et singularem erga se amorem perspexisset, teneri non potui quin tibi apertius illud idem his litteris declararem. ex omnibus enim, mi Dolabella, studiis in me et officiis quae summa sunt hoc scito mihi et amplissimum videri et gratissimum esse quod perfeceris ut Atticus intellegeret quantum ego te, quantum tu me amares.

\ciceroSection{3} quod reliquum est, Buthrotiam a et causam et civitatem, quamquam a te constituta est (beneficia autem nostra tueri solemus), tamen velim receptam in fidem tuam a meque etiam atque etiam tibi commendatam auctoritate et auxilio tuo tectam velis esse. satis erit in perpetuum Buthrotiis praesidi magnaque cura et sollicitudine Atticum et me liberaris, si hoc honoris mei causa susceperis ut eos semper a te defensos velis. quod ut facias te vehementer etiam atque etiam rogo.

\ciceroSection{4} his litteris scriptis me ad συντάξεισ dedi; quae quidem vereor ne miniata cerula tua pluribus locis notandae sint. ita sum μετέωροσ et magnis cogitationibus impeditus.
